
  The function
  \[
  f(x_1,x_2)=400-3x_1^2x_2^2
  \]is differentiable at
$(1,-2)^T$. Therefore the gradient is the direction of steepest
ascent. We have
\[
\nabla f(x_1,x_2)= \cvect{-6x_1 x_2^2 \\-6x_1^2 x_2},
\]
and
\[
\nabla f(1,-2)=\cvect{-24 \\12}.
\]
As $\|\nabla f(1,-2)\| = 12 \sqrt{5}$,
the normalized direction in which the depth increases as fast as
possible is then
\[
d = \cvect{-2/\sqrt{5},1/\sqrt{5}}.
\]
If she swims one meter along the direction, her new position is
\[
x^+ = \cvect{1 -2/\sqrt{5} \\ -2 + 1/\sqrt{5}},
\]
and the depth at this position is
\[
f(x^+) = 400 -3\frac{269-120\sqrt{5}}{25} \approx 399.92.
\]
If she swims a distance of 8 meters,
 her new position is
\[
x^{++} = \cvect{1 -16/\sqrt{5} \\ -2 + 8/\sqrt{5}},
\]
and the depth is
\[
f(x^{++}) \approx 117.
\]
As the depth at the current position is
\[
f(x) = 388,
\]
it illustrates that following an ascent direction increases the value
of the function, but up to a point. If the step along the direction is too
long, the value of the function may actually decrease. In this
example, the step should be lesser than 5.59 in order to obtain an
increase of the value of the function. 

